\chapter{Modul Pemindaian \& Operasional Absensi}

Modul Pemindaian adalah jantung operasional absensi harian EduConnect. Modul ini dirancang mampu memproses ratusan hingga ribuan pemindaian siswa di gerbang secara instan, baik melalui pemindai optik kamera (\textit{QR Scanner}) maupun pencatatan darurat langsung oleh guru piket (\textit{Manual Attendance}). Selain itu, modul \textbf{Rekap Kehadiran} memungkinkan Admin dan Wali Kelas untuk melihat, mengedit, dan menghapus rekaman absensi harian.

\section{Konsep Dasar \& Matriks Peran}
\begin{itemize}
    \item \textbf{Guru Piket / Staf (\textit{Scanner/Staff Role})}: Bertanggung jawab mengoperasikan perangkat (tablet/laptop) yang mengaktifkan kamera pemindai QR. Mereka hanya memiliki hak akses ke modul pemindai (\texttt{/scanner}) dan tidak dapat mengubah data \textit{Master}.
    \item \textbf{Admin, Super Admin, Teacher}: Semua peran ini juga dapat mengakses halaman Scanner, berguna saat guru piket berhalangan hadir.
    \item \textbf{Siswa}: Sebagai subjek yang menunjukkan kartu fisik atau tampilan QR dari perangkat \textit{smartphone} mereka.
\end{itemize}

\section{Skenario Dunia Nyata \& Alur Kerja}

Pukul 06:45 pagi, Guru Piket sudah menyalakan laptop dan membuka halaman \textbf{Scanner Absensi}. Fase 1: Siswa A melintas dan mendekatkan kartu QR-nya ke kamera. Dalam sepersekian detik, layar berkedip hijau dan menampilkan notifikasi sukses bertuliskan nama Siswa A beserta jam \textit{check-in}-nya. Antrean berjalan lancar. Fase 2: Siswa B datang menangis karena kartunya tertinggal di rumah. Guru Piket menenangkannya, lalu beralih ke tab \textbf{Absensi Manual}. Ia mengetikkan tiga huruf nama Siswa B, memilih nama yang muncul dari daftar sugesti, memilih status \textit{Hadir}, mengisi catatan \textit{``Kartu tertinggal''}, dan menekan tombol \textbf{Simpan Absensi}. Siswa B berhasil diabsen tanpa kartu.

\begin{figure}[H]
    \centering
    \includegraphics[width=0.9\textwidth]{placeholder_qr_scanner.png}
    \caption{Antarmuka halaman Scanner Absensi (/scanner) yang diakses oleh Guru Piket.}
    \label{figure:4.1}
\end{figure}

\section{Panduan Penggunaan (Step-by-Step)}

\subsection{Menjalankan Pemindai QR (Camera Scanner)}
\begin{enumerate}
    \item Pastikan perangkat Anda (Laptop/Tablet) terhubung ke internet dan memiliki kamera yang berfungsi.
    \item \textit{Login} ke sistem menggunakan akun berwewenang (Admin/Guru/Piket).
    \item Pada menu navigasi panel kiri, klik ikon/tautan \textbf{Scanner Absensi}. Perhatian: tautan ini akan membuka halaman baru (\textit{new tab}) karena menu ini memiliki atribut \textit{target=\_blank}.
    \item \textit{Browser} akan memunculkan \textit{pop-up} izin keamanan (\textit{``App wants to use your camera''}). Anda wajib mengeklik \textbf{Allow} (Izinkan) untuk mengaktifkan kamera.
    \item Kamera akan aktif. Arahkan kartu QR siswa ke dalam bingkai \textit{viewfinder} di tengah layar.
    \item Jika pemindaian sukses, layar akan berkedip \textbf{hijau} dan memunculkan notifikasi (\textit{toast}) berisi nama siswa dan jam \textit{check-in} mereka. Proses dapat dilanjutkan terus-menerus untuk siswa berikutnya tanpa perlu \textit{refresh}.
\end{enumerate}

\subsection{Melakukan Absensi Darurat (Manual Attendance)}
\begin{enumerate}
    \item Jika siswa tidak membawa kartu atau kamera rusak, pada halaman Scanner, cari dan klik tab atau bagian \textbf{Absensi Manual}.
    \item Anda akan melihat formulir pencarian dengan fitur \textit{autocomplete}. Klik pada kotak teks \textbf{Cari Nama atau NIS}.
    \item Mulai ketikkan minimal 3 huruf dari nama siswa. Sistem akan menampilkan daftar sugesti siswa yang cocok secara otomatis.
    \item Klik pada nama siswa yang tepat dari daftar sugesti tersebut.
    \item Pilih status absensi dari \textit{dropdown} \textbf{Status}. Pilihan yang tersedia adalah: \textbf{Hadir}, \textbf{Terlambat}, \textbf{Izin}, dan \textbf{Sakit}.
    \item Tambahkan keterangan tambahan pada isian teks \textbf{Catatan} jika diperlukan (contoh: \textit{``Hadir tanpa kartu''}).
    \item Selesaikan dengan menekan tombol biru \textbf{Simpan Absensi}.
\end{enumerate}

\subsection{Melihat \& Mengedit Rekap Kehadiran (/admin/attendance)}
\begin{enumerate}
    \item Pada panel navigasi kiri, klik menu \textbf{Rekap Kehadiran} di kelompok \textit{Kehadiran \& Laporan}.
    \item Halaman menampilkan tabel presensi harian dengan beberapa komponen filter di atas:
    \begin{itemize}
        \item \textbf{Filter Tanggal}: Pilih tanggal spesifik menggunakan \textit{date picker}.
        \item \textbf{Filter Kelas}: \textit{Dropdown} daftar kelas untuk menyaring per kelas.
        \item \textbf{Filter Status}: \textit{Dropdown} berisi pilihan \textbf{Semua Status}, Hadir, Terlambat, Izin, Sakit, Alpha.
        \item \textbf{Filter Baris}: \textit{Dropdown} untuk memilih jumlah baris per halaman (20, 50, 100, 200, atau \textbf{Tampilkan Semua}).
    \end{itemize}
    \item Klik tombol biru \textbf{Filter} untuk menerapkan pilihan. Tabel akan menampilkan kolom: \textbf{Tanggal}, \textbf{Siswa} (Nama + NIS), \textbf{Kelas}, \textbf{Jam Masuk}, \textbf{Jam Keluar}, \textbf{Status}, \textbf{Metode}, dan \textbf{Aksi}.
    \item Perhatikan indikator informasi di sudut kanan atas tabel (contoh: \textit{``Menampilkan 1 - 20 dari 840 data''}).
    \item Di bagian bawah tabel, terdapat navigasi halaman (\textit{pagination}). Gunakan tombol halaman untuk berpindah.
    \item Untuk mengedit satu rekaman absensi, klik tombol \textbf{Edit} di kolom Aksi pada baris yang diinginkan.
    \item Untuk menghapus rekaman absensi, klik tombol merah \textbf{Hapus}. Sistem akan menampilkan dialog konfirmasi browser sebelum mengeksekusi.
\end{enumerate}

\section{Penanganan Masalah (Edge Cases)}
\begin{itemize}
    \item \textbf{Pesan Error ``QR Code tidak dikenal''}: Token QR EduConnect di-\textit{generate} ulang setiap hari melalui mekanisme \textit{Cron Job} sistem. Kartu QR yang dicetak kemarin atau kartu bajakan tidak akan bisa digunakan pada hari ini. Siswa harus mengunduh tampilan QR versi hari ini melalui Portal Orang Tua atau meminta cetak ulang.
    \item \textbf{Siswa Terlambat Menimpa Status Alpha}: Jika \textit{Cron Job} sistem sudah terlanjur menandai siswa sebagai \textit{Alpha} karena belum hadir hingga jam yang ditentukan, lalu siswa tersebut datang terlambat dan men-\textit{scan} kartu QR-nya, sistem secara otomatis akan mengubah statusnya dari \textit{Alpha} menjadi \textit{Terlambat}. Sistem tidak akan salah menganggapnya sebagai \textit{Check-out} pulang.
    \item \textbf{Filter Rekap Kehadiran Tidak Menampilkan Data}: Pastikan tanggal yang dipilih pada filter sesuai dengan rentang tanggal data yang tersedia. Data absensi hanya tersedia untuk tanggal yang sudah dilewati dan bukan hari libur nasional (\textit{holiday}).
\end{itemize}
